\begin{codebox}
\codeline{\textcolor{cmtgray}{\#\ 幻觉控制三层架构}}
\codeline{\textcolor{cmtgray}{\#\ Three-Layer\ Hallucination\ Control\ Architecture}}
\codeline{\textcolor{cmtgray}{\#}}
\codeline{\textcolor{cmtgray}{\#\ 原则：不试图"训练掉"幻觉，而是用本体在系统层面拦截幻觉}}
\codeblank{}
\codeline{\textcolor{cmtgray}{\#\ ==============================}}
\codeline{\textcolor{cmtgray}{\#\ 1.\ 幻觉的根源}}
\codeline{\textcolor{cmtgray}{\#\ ==============================}}
\codeline{\textcolor{cmtgray}{\#\ LLM\ 生成的是"在语言上最可能"的续写，而非"经过核实"的事实：}}
\codeline{\textcolor{cmtgray}{\#\ \ \ 训练数据中没有\ Lathe\_003\ 的功率\ \(\rightarrow\)\ 模型按"车床通常15\textasciitilde{}30kW"编一个}}
\codeline{\textcolor{cmtgray}{\#\ \ \ 提问超出知识边界\ \(\rightarrow\)\ 模型倾向于给出流利但无依据的答案}}
\codeline{\textcolor{cmtgray}{\#\ 幻觉无法仅靠提示词消除——需要外部知识系统提供硬约束}}
\codeblank{}
\codeline{\textcolor{cmtgray}{\#\ ==============================}}
\codeline{\textcolor{cmtgray}{\#\ 2.\ 三层控制架构}}
\codeline{\textcolor{cmtgray}{\#\ ==============================}}
\codeline{\textcolor{cmtgray}{\#}}
\codeline{\textcolor{cmtgray}{\#\ \ \ 用户问题}}
\codeline{\textcolor{cmtgray}{\#\ \ \ \ \ │}}
\codeline{\textcolor{cmtgray}{\#\ \ \ [第一层\ 生成前]\ 约束注入（Prompt阶段）}}
\codeline{\textcolor{cmtgray}{\#\ \ \ \ \ │\ \ \ ·\ 把本体词汇表/类定义注入上下文："只能使用以下设备类型……"}}
\codeline{\textcolor{cmtgray}{\#\ \ \ \ \ │\ \ \ ·\ 把取值范围注入："功率单位kW，范围(0,10000]"}}
\codeline{\textcolor{cmtgray}{\#\ \ \ \ \ │\ \ \ ·\ 声明边界："知识库中没有的信息回答'无记录'"}}
\codeline{\textcolor{cmtgray}{\#\ \ \ \ \ │}}
\codeline{\textcolor{cmtgray}{\#\ \ \ [第二层\ 生成中]\ 工具调用（受控通道）}}
\codeline{\textcolor{cmtgray}{\#\ \ \ \ \ │\ \ \ ·\ LLM不直接回答事实，而是生成查询（SPARQL/函数调用）}}
\codeline{\textcolor{cmtgray}{\#\ \ \ \ \ │\ \ \ ·\ 事实只能来自知识图谱返回值（见\ text2sparql.txt）}}
\codeline{\textcolor{cmtgray}{\#\ \ \ \ \ │}}
\codeline{\textcolor{cmtgray}{\#\ \ \ [第三层\ 生成后]\ 事实校验（出口闸门）}}
\codeline{\textcolor{cmtgray}{\#\ \ \ \ \ │\ \ \ ·\ 实体校验：回答中的实体能否链接到本体个体？}}
\codeline{\textcolor{cmtgray}{\#\ \ \ \ \ │\ \ \ ·\ 断言校验：回答中的三元组是否存在于图谱/可由推理导出？}}
\codeline{\textcolor{cmtgray}{\#\ \ \ \ \ │\ \ \ ·\ 范围校验：数值是否通过SHACL形状检查？}}
\codeline{\textcolor{cmtgray}{\#\ \ \ \ \ │\ \ \ ·\ 一致性校验：是否违反不相交/函数性公理？}}
\codeline{\textcolor{cmtgray}{\#\ \ \ \ \ ▼}}
\codeline{\textcolor{cmtgray}{\#\ \ \ 通过\ \(\rightarrow\)\ 输出（附溯源）；不通过\ \(\rightarrow\)\ 拒答或降级为"无记录"}}
\codeblank{}
\codeline{\textcolor{cmtgray}{\#\ ==============================}}
\codeline{\textcolor{cmtgray}{\#\ 3.\ 制造场景演示}}
\codeline{\textcolor{cmtgray}{\#\ ==============================}}
\codeline{\textcolor{cmtgray}{\#\ 问："Lathe\_003\ 能加工钛合金吗？它功率多大？"}}
\codeblank{}
\codeline{\textcolor{cmtgray}{\#\ 无控制（直接生成）：}}
\codeline{\textcolor{cmtgray}{\#\ \ \ "Lathe\_003功率22kW，可以加工钛合金。"\ \ ←\ 两个事实都是编造}}
\codeblank{}
\codeline{\textcolor{cmtgray}{\#\ 三层控制下：}}
\codeline{\textcolor{cmtgray}{\#\ \ \ 第一层注入：本体中\ Equipment/Material\ 词汇\ +\ canProcess\ 定义}}
\codeline{\textcolor{cmtgray}{\#\ \ \ 第二层查询：}}
\codeline{\textcolor{cmtgray}{\#\ \ \ \ \ ASK\ \{\ mfg:Lathe\_003\ mfg:canProcess\ mfg:TitaniumAlloy\ \}\ \ \(\rightarrow\)\ false}}
\codeline{\textcolor{cmtgray}{\#\ \ \ \ \ SELECT\ ?p\ \{\ mfg:Lathe\_003\ mfg:hasPower\ ?p\ \}\ \ \ \ \ \ \ \ \ \ \ \ \ \(\rightarrow\)\ 15.5}}
\codeline{\textcolor{cmtgray}{\#\ \ \ 第三层校验后输出：}}
\codeline{\textcolor{cmtgray}{\#\ \ \ \ \ "功率15.5kW（来源：设备台账）。知识库中没有该设备可加工}}
\codeline{\textcolor{cmtgray}{\#\ \ \ \ \ \ 钛合金的能力记录（注意：'无记录'不等于'不能'，}}
\codeline{\textcolor{cmtgray}{\#\ \ \ \ \ \ 如需确认请联系工艺部门）。"}}
\codeblank{}
\codeline{\textcolor{cmtgray}{\#\ ==============================}}
\codeline{\textcolor{cmtgray}{\#\ 4.\ 类型混淆的公理拦截}}
\codeline{\textcolor{cmtgray}{\#\ ==============================}}
\codeline{\textcolor{cmtgray}{\#\ LLM常见错误：把概念归入错误类型（嫁接幻觉）}}
\codeline{\textcolor{cmtgray}{\#\ 防线：互不相交公理\ +\ 出口校验}}
\codeblank{}
\codeline{\textcolor{cmtgray}{\#\ 本体公理：}}
\codeline{\textcolor{cmtgray}{\#\ \ \ AllDisjointClasses(CNCLathe,\ CNCMillingMachine,\ Robot,\ AGV)}}
\codeline{\textcolor{cmtgray}{\#\ \ \ Equipment\ \(\sqcap\)\ Material\ \(\sqsubseteq\)\ \(\bot\)}}
\codeblank{}
\codeline{\textcolor{cmtgray}{\#\ 拦截示例：}}
\codeline{\textcolor{cmtgray}{\#\ \ \ LLM输出断言\ (Lathe\_003,\ rdf:type,\ Robot)}}
\codeline{\textcolor{cmtgray}{\#\ \ \ \(\rightarrow\)\ 与已有断言\ (Lathe\_003,\ rdf:type,\ CNCLathe)\ 在不相交公理下矛盾}}
\codeline{\textcolor{cmtgray}{\#\ \ \ \(\rightarrow\)\ 推理器报不一致\ \(\rightarrow\)\ 该输出被拒绝}}
\codeblank{}
\codeline{\textcolor{cmtgray}{\#\ 工程提示：不相交公理是性价比最高的防幻觉公理——}}
\codeline{\textcolor{cmtgray}{\#\ 建模时对每组同层兄弟类显式声明\ AllDisjointClasses}}
\codeblank{}
\codeline{\textcolor{cmtgray}{\#\ ==============================}}
\codeline{\textcolor{cmtgray}{\#\ 5.\ 溯源与审计（PROV思想）}}
\codeline{\textcolor{cmtgray}{\#\ ==============================}}
\codeline{\textcolor{cmtgray}{\#\ 每个对外输出附带证据链：}}
\codeline{\textcolor{cmtgray}{\#\ \ \ 结论：\ \ \ \ 功率15.5kW}}
\codeline{\textcolor{cmtgray}{\#\ \ \ 证据：\ \ \ \ (mfg:Lathe\_003,\ mfg:hasPower,\ 15.5)}}
\codeline{\textcolor{cmtgray}{\#\ \ \ 来源图：\ \ graph\ <http://example.org/graph/ledger>（台账，可信度高）}}
\codeline{\textcolor{cmtgray}{\#\ \ \ 查询：\ \ \ \ 所执行的SPARQL原文}}
\codeline{\textcolor{cmtgray}{\#\ \ \ 时间戳：\ \ 2026-06-11T10:30:00+08:00}}
\codeline{\textcolor{cmtgray}{\#\ 价值：事后可审计"AI为什么这么说"；}}
\codeline{\textcolor{cmtgray}{\#\ 高合规行业（航空DO-178C、汽车ISO\ 26262、医药GxP）可直接复用该审计链}}
\codeblank{}
\codeline{\textcolor{cmtgray}{\#\ ==============================}}
\codeline{\textcolor{cmtgray}{\#\ 6.\ 评测指标}}
\codeline{\textcolor{cmtgray}{\#\ ==============================}}
\codeline{\textcolor{cmtgray}{\#\ 上线前用标准问题集评测：}}
\codeline{\textcolor{cmtgray}{\#\ \ \ 事实准确率\ \ \ ——\ 回答中可校验断言的正确比例（目标\ \(\ge\)\ 99\%）}}
\codeline{\textcolor{cmtgray}{\#\ \ \ 拒答正确率\ \ \ ——\ 知识库无记录时正确说"无记录"的比例}}
\codeline{\textcolor{cmtgray}{\#\ \ \ 幻觉漏过率\ \ \ ——\ 编造内容通过三层闸门的比例（目标\ \(\rightarrow\)\ 0）}}
\codeline{\textcolor{cmtgray}{\#\ 经验：第二层（工具调用）贡献最大；三层叠加后漏过率可降1\textasciitilde{}2个数量级}}
\end{codebox}
